% Generated draft: requires mathematical and source review.
\documentclass[11pt]{article}
\usepackage[margin=25mm]{geometry}
\usepackage{fontspec}
\setmainfont{Latin Modern Roman}
\usepackage{amsmath,amssymb,mathtools,unicode-math}
\setmathfont{Latin Modern Math}
\usepackage{enumitem,needspace,xurl,hyperref}
\setlength{\parindent}{0pt}
\setlength{\parskip}{5pt}
\setlength{\emergencystretch}{3em}
\newcommand{\N}{\mathbb{N}}
\newcommand{\Z}{\mathbb{Z}}
\newcommand{\Q}{\mathbb{Q}}
\newcommand{\R}{\mathbb{R}}
\newcommand{\C}{\mathbb{C}}
\newcommand{\F}{\mathbb{F}}
\newcommand{\A}{\mathbb{A}}
\newcommand{\PP}{\mathbb{P}}
\newcommand{\E}{\mathbb{E}}
\newcommand{\eps}{\varepsilon}
\newcommand{\dd}{\,\mathrm{d}}
\begin{document}
\section*{UnsolvedMath ID: 30001267}
\textbf{Source problem code:} catalogue entry
\subsection{Definitions and mathematical statement}
The supplied record contains four questions. The following gives a precise conventional interpretation while retaining the unresolved terminology and domain qualifications.

For a finitely generated group \(G\) with finite symmetric generating set \(S\), let \(d_S(g,h)=|g^{-1}h|_S\) be its word metric. For a metric group \((G,d)\), define its non-equivariant Hilbert-space compression exponent by
\[
\alpha_2(G,d)=\sup\left\{\alpha\in[0,1]:\begin{array}{l}
\text{there exist a real Hilbert space }H,\ f:G\to H,\\
 c,L>0,\ C\geq0\text{ such that, for all }g,h\in G,\\
 c\,d(g,h)^\alpha-C\leq\|f(g)-f(h)\|\leq Ld(g,h)+C
\end{array}\right\}.
\]
The lower bound is required only when \(d(g,h)\geq1\), avoiding any convention about \(0^0\). For word metrics the exponent is independent of the finite generating set. No equivariance of \(f\) is required. The supplied record does not explicitly distinguish this exponent from equivariant compression.

Here a \emph{uniform embedding}, in the coarse-geometric sense, of \((G,d)\) into a Banach space \(X\) means a map \(f:G\to X\) and nondecreasing functions \(\rho_-,\rho_+:[0,\infty)\to[0,\infty)\), with \(\rho_-(t)\to\infty\) as \(t\to\infty\), such that
\[
\rho_-(d(g,h))\leq\|f(g)-f(h)\|\leq\rho_+(d(g,h))
\qquad(g,h\in G).
\]
A group is uniformly embeddable into Hilbert space if such an embedding exists with \(X\) a Hilbert space. This interpretation of uniform embeddability requires confirmation from the report.

A discrete group \(G\) is \emph{amenable} if there is a positive linear functional \(m:\ell^\infty(G)\to\R\) on bounded real-valued functions such that \(m(1)=1\) and \(m(u\circ L_g)=m(u)\) for every \(g\in G\), where \(L_g(h)=gh\).

Under the Cowling--Haagerup interpretation, a discrete group \(G\) is \emph{weakly amenable} if there is a net of finitely supported functions \(\varphi_i:G\to\C\) converging pointwise to \(1\), whose associated Schur multipliers
\[
M_{\varphi_i}\bigl((T_{g,h})_{g,h\in G}\bigr)
 =\bigl(\varphi_i(g^{-1}h)T_{g,h}\bigr)_{g,h\in G}
\]
on \(B(\ell^2(G))\) have uniformly bounded completely bounded norms. Here \(\|M\|_{\mathrm{cb}}=\sup_{n\geq1}\|\operatorname{id}_{M_n}\otimes M\|\). The record does not specify the intended meaning of weak amenability, so this definition is conditional.

A Banach space \(X\) is \emph{uniformly convex} if for every \(\varepsilon\in(0,2]\) there is \(\delta>0\) such that
\[
\|x\|,\|y\|\leq1,\quad\|x-y\|\geq\varepsilon
\quad\Longrightarrow\quad
\left\|\frac{x+y}{2}\right\|\leq1-\delta.
\]

Let \(\mathfrak G\) denote the group called Grigorchuk's group of intermediate growth in the supplied record. The usual identification with the first Grigorchuk group is an interpretation requiring source confirmation; the record supplies no defining action or presentation. Let \(\mathcal C\) denote the intended class of discrete groups equipped with metrics. The source does not determine \(\mathcal C\) or its metric convention for Questions 2 and 3. Taking finitely generated groups with word metrics gives a precise restricted reading, but must not be silently substituted for the full intended scope.

Subject to these explicit qualifications, the four questions are:
\begin{enumerate}
\item Determine \(\alpha_2(\mathfrak G,d_S)\), where \(S\) is a finite symmetric generating set of the intended Grigorchuk group.
\item Does there exist \((G,d)\in\mathcal C\) such that \(G\) is amenable and \(\alpha_2(G,d)<\tfrac12\)?
\item Is it true that every \((G,d)\in\mathcal C\) which admits a uniform embedding into a Hilbert space has weakly amenable underlying group \(G\)?
\item Does there exist a finitely generated group \(G\), equipped with a word metric, such that for every uniformly convex Banach space \(X\), there is no uniform embedding \(G\to X\)?
\end{enumerate}
These are statement-only questions. The unspecified class, terminology, and group identification prevent this conditional formulation from being treated as a fully verified transcription of the report.
\subsection{Short English statement}
The record asks for the Hilbert-space compression of Grigorchuk's intermediate-growth group, whether amenable groups can have compression below \(1/2\), whether Hilbert-space uniform embeddability implies weak amenability, and whether some finitely generated group fails to embed uniformly into every uniformly convex Banach space. The meanings of compression, uniform embedding, and weak amenability, and the domains of the middle two questions, require source confirmation.
\subsection{Sources}
\begin{itemize}
\item Local UnsolvedMath registry, permanent ID \(30001267\), source problem code \texttt{catalogue entry}, title \emph{Compression, Amenability, and Uniform Embeddability of Groups}; collection: Oberwolfach Reports Open Problems; proposed year: 2009. No individual proposer is supplied.
\item Supplied source citation: \emph{Manifold Perspectives} (2009), Oberwolfach Reports, pp.\ 1487--1546, DOI \texttt{10.4171/owr/2009/27}. Supplied source URL and literature-source label ``original OWR report'': \url{https://doi.org/10.4171/owr/2009/27}. No page or passage locating these four questions within the report is supplied.
\item Supplied literature source, ``compression constructions'': \url{https://arxiv.org/abs/math/0612378}. The record supplies no author, title, or precise result for this reference.
\item Supplied literature source, ``report open-problem list'': \url{https://maths.ed.ac.uk/~v1ranick/papers/OWR_2009_27.pdf}.
\item Literature assessment dated exactly \texttt{2026-08-21}: the supplied curation labels the problem open and states that the report lists four open questions, while later work gives related compression and non-embeddability results without settling the full list. This is supplied dated curation, not an independently verified literature assessment. The background expressly requests independent verification of time-sensitive status and cited claims; no source has been visited for this draft.
\end{itemize}
\end{document}
